\section{Adding the prime costs}\label{growth:section}

Fix \(M=400\), \(\alpha=3/40\), and \(\lambda=37/40\).
For sufficiently large \(K=40n\), the preceding bounds give integers
\begin{equation}\label{growth:Lp}
 L_p=
 \begin{cases}
 -6h\lfloor\log_p(5K)\rfloor-hv_p(24),&p\le K/400,\\
 v_p(S_K)+\lfloor2W_p^{\rm in}\rfloor,&K/400<p\le K/3,\\
 v_p(S_K)+\lfloor2W_p^{\rm out}-r_p\rfloor,&K/3<p\le2h,
 \end{cases}
\end{equation}
with \(L_p\le v_p^{\mathrm G}(F_K)\).
For \(p>2h\), the same outer calculation has no large-pole rows or
rank loss, and the factorial scalar is a unit; hence
\(v_p^{\mathrm G}(F_K)\ge0\).
This is the normalizer in Firsching's revised argument
\cite[\texttt{Norm.lean}]{firsching}.

Set
\begin{equation}\label{growth:normalizer}
 m_K=\prod_{p\le2h}p^{-L_p}.
\end{equation}
Every coefficient of \(m_KF_K\) has nonnegative valuation at every
prime, so its reduced denominator is one. Consequently
\begin{equation}\label{growth:integrality}
 m_KF_K\in\Z[X],\qquad
 \log m_K=\sum_{p\le2h}(-L_p)\log p.
\end{equation}
The multiplier is positive and rational. Positive \(L_p\)'s allow
common factors to be divided out and contribute \emph{negative} costs;
discarding those savings would spoil the estimate.

\subsection{From local counts to a fixed budget}

The local formulas still depend on \(K\) and \(p\).
Writing \(x=K/p\) converts them into bounded, piecewise continuous
profiles on fixed intervals. First the residue classes \(c/p\)
are averaged; then the resulting costs are summed over primes.
The profiles retain both the basis weights and the factorial
valuation \(v_p(S_K)\), computed by Legendre's formula.
Their explicit definitions are in companion Sections 7.4 and 7.7.
For each range above the small-prime cutoff, the corresponding profile
\(E\) satisfies
\[
 -L_p\le pE(K/p)+O_{400}(1).
\]
This converts the local divisibility bound into a cost that can be summed.

The prime number theorem,
\(\vartheta(T)=\sum_{p\le T}\log p=T+o(T)\), gives
\begin{equation}\label{growth:prime-sum}
 \frac1{K^2}\sum_{K/b<p\le K/a}p\log p\,E(K/p)
 \longrightarrow\int_a^b\frac{E(x)}{x^3}\,dx
\end{equation}
for a bounded piecewise continuous \(E\) with finitely many
discontinuities on a fixed \(0<a<b\).
To see the factor \(x^{-3}\), the prime number theorem replaces
\(K^{-1}\sum(\log p)\delta_{p/K}\) by length measure; apply this to
\(tE(1/t)\), then substitute \(x=1/t\).
An \(O_{400}(1)\) error at each prime contributes only
\(O_{400}(\vartheta(2h))=O(K)=o(K^2)\).
This is the prime-summation step of Fauzan and Firsching
\cite[Section 5]{fauzan}\cite[\texttt{PrimeSumIntegral.lean}]{firsching}.

There are four contributions. Small primes use the coarse bound:
apart from \(O(\sqrt K)\) primes, the logarithmic floor equals one,
giving \(6\lambda/400=0.013875\).
For \(3\le x\le20\), the inner profile has 125 rational pieces,
integrated exactly. For \(20\le x\le400\), a quadratic bound followed
by two integrations by parts gives a negative contribution.
The successive primitives are periodic Bernoulli polynomials: they
remove the oscillating fractional-part counting error.
The outer profile has ten merged rational pieces and includes the
full rank correction. Companion Sections 7.3--7.8 give these calculations;
the supplied exact-arithmetic checks reproduce their values.

\begin{center}
\begin{tabular}{l r}
\toprule
Prime range & Upper bound for its \(K^2\) coefficient\\
\midrule
\(p\le K/400\) & \(0.013875\)\\
\(K/400<p\le K/20\) & \(-0.05346\)\\
\(K/20<p\le K/3\) & \(0.04279\)\\
\(K/3<p\le2h\) & \(1.34481\)\\
\bottomrule
\end{tabular}
\end{center}
These finite decimals are rational upper bounds, rounded upwards
from the exact values. Their sum is \(1.348015<1.35\).
The exact calculation gives the slightly stronger bound
\[
 \limsup_{\substack{K\to\infty\\40\mid K}}
 \frac{\log m_K}{K^2}
 \le B=1.348003127429204\ldots<1.35.
\]
The exact fraction defining \(B\) is recorded with the final comparison
below and in the companion's Section 7.8.
Because the gap below \(1.35\) is strict, all sufficiently large
admissible \(K\) satisfy
\begin{equation}\label{growth:final}
 \boxed{m_K>0,\qquad m_KF_K\in\Z[X],\qquad
        \log m_K\le1.35K^2.}
\end{equation}
This is the arithmetic budget to compare with the real decay.
